\documentclass[10pt,a4paper]{amsart}
\usepackage{amsmath,amssymb,amsthm,graphicx}
\usepackage{verbatim}
\newenvironment{reference}{\comment}{\endcomment}
%\newenvironment{reference}{}{}
\newenvironment{slogan}{\comment}{\endcomment}
\newenvironment{history}{\comment}{\endcomment}

% For commutative diagrams we use Xy-pic
\usepackage[all]{xy}

% We use 2cell for 2-commutative diagrams.
\xyoption{2cell}
% Theorem environments.
%
\theoremstyle{plain}
\newtheorem{theorem}[subsection]{Theorem}
\newtheorem{proposition}[subsection]{Proposition}
\newtheorem{lemma}[subsection]{Lemma}

\theoremstyle{definition}
\newtheorem{definition}[subsection]{Definition}
\newtheorem{example}[subsection]{Example}
\newtheorem{exercise}[subsection]{Exercise}
\newtheorem{situation}[subsection]{Situation}

\theoremstyle{remark}
\newtheorem{remark}[subsection]{Remark}
\newtheorem{remarks}[subsection]{Remarks}

\numberwithin{equation}{subsection}

% Macros
%
\def\lim{\mathop{\mathrm{lim}}\nolimits}
\def\colim{\mathop{\mathrm{colim}}\nolimits}
\def\Spec{\mathop{\mathrm{Spec}}}
\def\Hom{\mathop{\mathrm{Hom}}\nolimits}
\def\Ext{\mathop{\mathrm{Ext}}\nolimits}
\def\SheafHom{\mathop{\mathcal{H}\!\mathit{om}}\nolimits}
\def\SheafExt{\mathop{\mathcal{E}\!\mathit{xt}}\nolimits}
\def\Sch{\mathit{Sch}}
\def\Mor{\mathop{\mathrm{Mor}}\nolimits}
\def\Ob{\mathop{\mathrm{Ob}}\nolimits}
\def\Sh{\mathop{\mathit{Sh}}\nolimits}
\def\NL{\mathop{N\!L}\nolimits}
\def\CH{\mathop{\mathrm{CH}}\nolimits}
\def\proetale{{pro\text{-}\acute{e}tale}}
\def\etale{{\acute{e}tale}}
\def\QCoh{\mathit{QCoh}}
\def\Ker{\mathop{\mathrm{Ker}}}
\def\Im{\mathop{\mathrm{Im}}}
\def\Coker{\mathop{\mathrm{Coker}}}
\def\Coim{\mathop{\mathrm{Coim}}}

% Boxtimes
%
\DeclareMathSymbol{\boxtimes}{\mathbin}{AMSa}{"02}

%
% Macros for moduli stacks/spaces
%
\def\QCohstack{\mathcal{QC}\!\mathit{oh}}
\def\Cohstack{\mathcal{C}\!\mathit{oh}}
\def\Spacesstack{\mathcal{S}\!\mathit{paces}}
\def\Quotfunctor{\mathrm{Quot}}
\def\Hilbfunctor{\mathrm{Hilb}}
\def\Curvesstack{\mathcal{C}\!\mathit{urves}}
\def\Polarizedstack{\mathcal{P}\!\mathit{olarized}}
\def\Complexesstack{\mathcal{C}\!\mathit{omplexes}}
% \Pic is the operator that assigns to X its picard group, usage \Pic(X)
% \Picardstack_{X/B} denotes the Picard stack of X over B
% \Picardfunctor_{X/B} denotes the Picard functor of X over B
\def\Pic{\mathop{\mathrm{Pic}}\nolimits}
\def\Picardstack{\mathcal{P}\!\mathit{ic}}
\def\Picardfunctor{\mathrm{Pic}}
\def\Deformationcategory{\mathcal{D}\!\mathit{ef}}
\title{Tautologically equivalent families give the same sheaf condition}
\author{The Stacks Project Authors}
\date{stacks-project:sites.tex:local-snapshot-20261012; GNU FDL 1.2 or later}
\begin{document}
\maketitle
\noindent\textit{Editorial note.} All delivered mathematical text is the author's own, copied from the pinned Stacks Project chapter file sites-cohomology.tex; no statement or proof step was rewritten, reordered or removed.\par\medskip
\noindent\textit{Editorial note.} The proof is the author's complete argument; every step the author writes out is inside the counted span and no closing sentence is a placeholder.\par\medskip
\noindent\textit{Editorial note.} Every labelled reference occurring inside the counted spans is carried in assembly.external\_references as a non-load-bearing navigation pointer (same-chapter ones at their definition site, cross-chapter ones at the invocation inside the counted span), so no load-bearing mathematics is replaced by a stub. No bibliography citation occurs inside the counted spans.\par\medskip
\noindent\textit{Editorial note.} The checked chapter list for same-name collisions is the set of stacks chapters staged in this repository: chow, cohomology, divisors, etale-cohomology, flat, more-algebra, more-morphisms, morphisms, spaces-more-morphisms, varieties together with this card's chapter sites-cohomology.\par\medskip
\noindent\textit{Editorial note.} The preamble is the chapter's own retained minimal preamble (sources/stacks-sites-cohomology/retained-preamble.tex), a verbatim excerpt of the native preamble: the Xy-pic setup and the native theorem declarations with the mathematical macros.\par\medskip
\noindent\textit{Allowed standard prerequisites (editorial).}\begin{itemize}
\item Tautological equivalence of families of morphisms. Same chapter, the definition of tautologically equivalent families with a fixed target. Supplies the index map and the isomorphisms over the target whose existence the statement assumes.
\end{itemize}\medskip
\begin{lemma}
\label{lemma-tautological-same-sheaf}
Let $\mathcal{C}$ be a category. Let
$\mathcal{U} = \{\varphi_i : U_i \to U\}_{i\in I}$, and
$\mathcal{V} = \{\psi_j : V_j \to U\}_{j\in J}$ be two families of morphisms
with the same fixed target. Assume that the fibre products
$U_i \times_U U_{i'}$ and $V_j \times_U V_{j'}$ exist.
If $\mathcal{U}$ and $\mathcal{V}$ are
tautologically equivalent, then for any presheaf $\mathcal{F}$ on
$\mathcal{C}$ the sheaf condition for $\mathcal{F}$ with respect to
$\mathcal{U}$ is equivalent to the sheaf condition for $\mathcal{F}$
with respect to $\mathcal{V}$.
\end{lemma}

\begin{proof}
First, note that if $\varphi : A \to B$ is an isomorphism in the
category $\mathcal{C}$, then $\varphi^* : \mathcal{F}(B) \to \mathcal{F}(A)$
is an isomorphism. Let $\beta : J \to I$ be a map and let
$\chi_j : V_j \to U_{\beta(j)}$ be isomorphisms over $U$ which
are assumed to exist by hypothesis. Let us show that the sheaf
condition for $\mathcal{V}$ implies the sheaf condition for $\mathcal{U}$.
Suppose given sections $s_i \in \mathcal{F}(U_i)$ such that
$$
s_i|_{U_i \times_U U_{i'}} = s_{i'}|_{U_i \times_U U_{i'}}
$$
in $\mathcal{F}(U_i \times_U U_{i'})$ for all pairs $(i, i') \in I \times I$.
Then we can define $s_j = \chi_j^*s_{\beta(j)}$. For any pair
$(j, j') \in J \times J$ the morphism
$\chi_j \times_{\text{id}_U} \chi_{j'} : V_j \times_U V_{j'} \to
U_{\beta(j)} \times_U U_{\beta(j')}$ is an isomorphism as well.
Hence by transport of structure we see that
$$
s_j|_{V_j \times_U V_{j'}} = s_{j'}|_{V_j \times_U V_{j'}}
$$
as well. The sheaf condition w.r.t.\ $\mathcal{V}$ implies there
exists a unique $s$ such that $s|_{V_j} = s_j$ for all $j \in J$.
By the first remark of the proof this implies that $s|_{U_i} = s_i$
for all $i \in \Im(\beta)$ as well. Suppose that $i \in I$,
$i \not \in \Im(\beta)$. For such an $i$ we have isomorphisms
$U_i \to V_{\alpha(i)} \to U_{\beta(\alpha(i))}$ over $U$. This gives a
morphism $U_i \to U_i \times_U U_{\beta(\alpha(i))}$ which is a
section of the projection. Because $s_i$ and $s_{\beta(\alpha(i))}$
restrict to the same element on the fibre product we conclude that
$s_{\beta(\alpha(i))}$ pulls back to $s_i$ via $U_i \to U_{\beta(\alpha(i))}$.
Thus we see that also $s_i = s|_{U_i}$ as desired.
\end{proof}

\end{document}
